\section{Discussion and recommendations}
\label{sec:discussion}

\paragraph{Answers to the research questions.}
\RQ{1}: the Runge--Kutta methods, fixed-step or adaptive, drift linearly in the mass shell and
quadratically in phase. The symplectic methods keep the mass-shell error bounded at
$\order(h^p)$ and the phase error linear, and once their truncation error falls below
round-off, compensated summation gives Brouwer's $N^{1/2}$.
\RQ{2}: in Hamilton's form $E$ and $L_z$ are exact for every method; in the second-order form
they drift for every non-geometric method. Formulation (b) is up to 440 times more accurate
for rays that go far out. Symmetry alone gives Gauss--Legendre in (a) the same bounded errors
as in (b). The one case that favours (a) is the near-polar orbit, where the spurious node
precession is about 100 times smaller.
\RQ{3}: DOP853 is the cheapest method over a few periods at every accuracy and over 1000
periods for modest accuracy; for a phase error of $10^{-8}$ over 1000 periods only GL3
delivers. Tao4 is never competitive.
\RQ{4}: on the photon sphere accuracy buys survival logarithmically, $\ln 10/2\pi$ orbits per
decade, up to about six orbits in double precision for every method. At the ISCO the time to
plunge grows algebraically, as $|\varepsilon|^{-1/4}$, and an error of the stabilizing sign
prevents the plunge.

\begin{table*}[t]
\caption{\label{tab:recommendations}Recommendations (T7), with the evidence they rest on.}
\centering\small
\begin{tabular}{>{\raggedright\arraybackslash}p{0.24\textwidth}>{\raggedright\arraybackslash}p{0.30\textwidth}>{\raggedright\arraybackslash}p{0.38\textwidth}}
\toprule
Task & Recommendation & Evidence \\
\midrule
Ray tracing, short integrations (deflection, a few orbits)
  & DOP853 in formulation (b)
  & Cheapest at every accuracy (\tab{cost}); exact $E$, $L_z$ and up to 440 times more
    accurate than (a) near $b_c$ (\fig{windings}). \\ \addlinespace
Long orbit evolution where the phase matters ($10^3$ periods and more)
  & GL3 in (b), iterated to stagnation (or to $10^{-12}$), with compensated summation
  & The only method to reach a phase error of $10^{-8}$ over 1000 periods (\tab{long-cost});
    bounded invariants (\fig{energy}); Brouwer's law (\fig{roundoff}). \\ \addlinespace
Long runs at modest accuracy
  & DOP853 is cheaper, but its invariants drift linearly
  & \tab{long-cost}, \fig{energy}, \fig{phase}. \\ \addlinespace
Studies that need bounded invariants (Poincar\'e sections, statistics of chaotic orbits)
  & A Gauss method with fixed step, in either formulation
  & Bounded $H$, $L^2$ in (a) and (b) (\fig{symmetric}, \fig{inclined}). \\ \addlinespace
Unstable orbits (photon sphere)
  & Any accurate method with compensated summation; expect at most $\sim$6 orbits
  & \fig{photon}, \tab{photon}; plain summation fakes survival (\fig{roundoff}). \\ \addlinespace
Explicit symplectic integration of a general metric
  & Tao's method only with $\omega$ tuned for the problem at hand
  & Fails on the ISCO and on inclined orbits with $\omega$ tuned on a stable orbit
    (Sec.~\ref{sec:method-details}). \\
\bottomrule
\end{tabular}
\end{table*}

\paragraph{Recommendations.} \tab{recommendations} turns these answers into advice (\RQ{5}).
Two conclusions go beyond the hypotheses we started from. First, the cost advantage of the
adaptive pair is large enough that geometric integration pays off only when the integration
is long \emph{and} the phase matters. Over $10^3$ periods a mass-shell error of $10^{-10}$
is cheaper to buy with a tight tolerance than with a symplectic method, although the drift
of DOP853 will overtake the bounded error of GL3 on longer runs. Second, the round-off details matter: compensated
summation turns a quadratic phase drift into a few ulp, and its absence can make an
integration look more stable than the physics allows.

\paragraph{Limitations.} The study uses fixed steps for the geometric methods; variable steps
destroy their long-time properties unless the time is transformed~\cite{calvo1993,gni2006},
and Ref.~\cite{seyrich2012} shows one way to combine them. Only Schwarzschild is covered,
where $t$ and $\phi$ are cyclic and the problem is integrable; in nonintegrable systems
(charged or spinning particles, perturbed metrics) the advantage of symplectic methods should
be larger. Wall times compare our own implementations, whose generic step code costs 2--5
times the bare vector field, so they are not a statement about the algorithms' intrinsic
cost. The explicit symplectic splittings of Wang et al.~\cite{wang2021} and Taylor
methods~\cite{biscani2021} were not tested.

\paragraph{Extension to Kerr.} The metric interface takes the full inverse metric and its
derivatives, and none of the integrators or formulations knows the metric. Adding Kerr in
Boyer--Lindquist coordinates is a new metric class, analytic references (Carter's constant,
spherical photon orbits) and new test cases, with no change to the integrator code. The
inclined orbits of TC6 already exercise the code path that will monitor Carter's constant.
